\documentclass[oneside]{scrbook}

% \usepackage[utf8]{inputenc}
\usepackage[T1]{fontenc}
\usepackage{version}
% \usepackage[latin1]{inputenc}
\usepackage[italian]{babel}
\usepackage{algpseudocode}
\usepackage{amsmath}
\usepackage{booktabs}
\usepackage{color}
\usepackage{amsfonts}
\usepackage{url}
\usepackage{algorithm}
\usepackage{chngcntr}
\usepackage[answerdelayed]{exercise}
\usepackage{graphicx}
\usepackage{caption}
\usepackage{listings}
\usepackage{tikz}
\usepackage{pgf}
\usepackage{colortbl}
\usepackage{listingsutf8}
\usepackage[pdftitle = {Dispense di esercizi}, pdfauthor = {Gianmarco Tedeschi e Alberto Marchesi}, pdfsubject = {Fondamenti di Informatica ($082746$ - A.A.~$2025-2026$)}, pdfcreator = {Gianmarco Tedeschi e Alberto Marchesi}, pdfproducer = {-}]{hyperref}
% \usepackage[automark,nouppercase]{scrpage2}
\usepackage{mdframed}
\usepackage{scrlayer-scrpage} % Necessary package for page styles

%\pagestyle{scrheadings}

%\usepackage{titlesec}
%\newcommand{\sectionbreak}{\clearpage}
\counterwithin{Exercise}{chapter}
\counterwithin{Answer}{chapter}

\algrenewcommand\algorithmicthen{\textbf{Allora}}
\algrenewcommand\algorithmicif{\textbf{Se}}
\algrenewcommand\algorithmicelse{\textbf{Altrimenti}}
\algrenewcommand\algorithmicend{\textbf{Chiudi}}
\algrenewcommand\algorithmicwhile{\textbf{Finch\'e}}
\algrenewcommand\algorithmicdo{\textbf{Esegui}}
\algrenewcommand{\algorithmiccomment}[1]{\hskip3em\texttt{//#1}}

\setlength{\parindent}{0in}
\setlength{\parskip}{10pt}

\newmdenv[%
        outerlinewidth=2,%
        roundcorner=10pt,%
        leftmargin=0,%
        rightmargin=0,%
        outerlinecolor=black,
        innertopmargin=\topskip,%
        splittopskip=\topskip,%
]{algo}


\ohead{\headmark}
\ihead{}
\chead{}
\ifoot{A.A.~$2025-2026$}
\cfoot{Fondamenti di Informatica (082746) - G. Tedeschi}
\ofoot{Pag.~\pagemark}
\setheadsepline{.4pt}
\setfootsepline{.4pt}

\definecolor{lbcolor}{rgb}{0.9,0.9,0.9}

\lstdefinestyle{cstyle}{%
  inputencoding=utf8/latin1,
  breaklines=true,
  breakatwhitespace=false,
  basicstyle=\small\ttfamily,
  language=c,
  frame=single,
  showstringspaces=false
}

\lstdefinestyle{cstylesmall}{%
  inputencoding=utf8/latin1,
  breaklines=true,
  breakatwhitespace=false,
  basicstyle=\scriptsize\ttfamily,
  language=c,
  frame=single,
  showstringspaces=false
}

\lstdefinestyle{console}{%
  inputencoding=utf8/latin1,
  breaklines=true,
  breakatwhitespace=false,
  basicstyle=\scriptsize\ttfamily,
  showstringspaces=false,
  backgroundcolor=\color{lbcolor}
}


\lstnewenvironment{clang}{\lstset{style=cstyle}}{}
\lstnewenvironment{clangsmall}{\lstset{style=cstylesmall}}{}

\addtocounter{chapter}{8}
\begin{document}

\chapter{File, Argomenti da Riga di Comando e Scope}

Questa dispensa approfondisce tre aspetti fondamentali del linguaggio C:
\textbf{gestione dei file}, \textbf{acquisizione degli argomenti da riga di comando} e \textbf{regole dello scope delle variabili}.
Si tratta di strumenti fondamentali per rendere i programmi più flessibili e in grado di interagire con l'esterno.

\section{I File in C}

Un \textbf{file} è una struttura dati persistente memorizzata su memoria di massa.
A differenza delle variabili in memoria, i file mantengono i dati anche dopo la fine dell'esecuzione del programma.

Per manipolare un file in C si utilizza una variabile puntatore di tipo \texttt{FILE*} che rappresenta il \textbf{descrittore del file}.

\subsection{Apertura di un File}

\begin{lstlisting}[style=cstyle]
FILE *fp;
fp = fopen("dati.txt", "r");  // apertura in lettura
\end{lstlisting}

La funzione \texttt{fopen} restituisce:
\begin{itemize}
    \item un puntatore a \texttt{FILE} (inizializzato all'inizio del file) se l'apertura ha successo.
    \item \texttt{NULL} in caso di errore.
\end{itemize}

\paragraph{Modalità di apertura più comuni}
\begin{itemize}
    \item \texttt{"r"} lettura (testo)
    \item \texttt{"w"} scrittura (crea o sovrascrive)
    \item \texttt{"a"} append
    \item \texttt{"rb"}, \texttt{"wb"}, \texttt{"ab"} versioni binarie
\end{itemize}

Dopo l'apertura, il programma opera sul file utilizzando la variabile che lo rappresenta.
Al termine, la corrispondendenza dovrà essere eliminata: operazione di \emph{chiusura} del file.

\subsection{Chiusura di un File}

\begin{lstlisting}[style=cstyle]
fclose(fp);
\end{lstlisting}

Il valore restituto è un intero:
\begin{itemize}
    \item 0 se la chiusura ha successo.
    \item EOF (valore intero negativo) in caso di errore.
\end{itemize}

Inoltre, l'utilizzo della funzione \texttt{fclose} assicura che tutti i dati eventualmente rimasti in buffer vengano scritti su disco.

\subsection{Lettura e Scrittura di un File}
Un file di testo è un file che contiene sequenze di caratteri. Per questo motivo un file
possiede delle caratterische specifiche che lo distinguono da un file binario:

\begin{itemize}
    \item esiste un concetto di riga e di fine riga ('\textbackslash n').
    \item i dati sono rappresentati come sequenze di caratteri ASCII.
    \item la sequenza di caratteri è terminata dal carattere speciale EOF.
\end{itemize}

I canali di I/O standard non sono altro che file di testo speciali:
\begin{itemize}
    \item \texttt{stdin} (standard input): tastiera
    \item \texttt{stdout} (standard output): schermo
    \item \texttt{stderr} (standard error): schermo (per messaggi di errore)
\end{itemize}

\paragraph{Lettura e Scrittura Formattata}

Per leggere o scrivere in un file di testo in maniera formattata ed effettuando
le conversioni richieste dove necessario si usano \texttt{fscanf} e \texttt{fprintf}.

\begin{lstlisting}[style=cstyle]
fprintf(fp, "%d %c", numero, carattere);
fscanf(fp, "%d %c", &numero, &carattere);
\end{lstlisting}

La lettura restituisce un intero che indica il numero di \emph{campi} letti con successo.
La scrittura restituisce il numero di \emph{caratteri} scritti con successo.

\textbf{Lettura Carattere per Carattere}

\begin{lstlisting}[style=cstyle]
int c;
c = fgetc(fp);
while (c != EOF) {
    putchar(c);
    c = fgetc(fp);
}
\end{lstlisting}

\textbf{Lettura di Linee di Testo}

\begin{lstlisting}[style=cstyle]
char buffer[100];
while (fgets(buffer, 100, fp) != NULL) {
    fputs(buffer, stdout);
}
\end{lstlisting}

\subsection{File Binari: fread e fwrite}
Un file binario è una sequenza di byte senza formattazione specifica.
Può essere usato per archiviare qualunque tipo di informazione. In questo caso input e output
avvengono sotto forma di una sequenza di byte.

Rispetto ai file di testo, con i file binari la fine del file non è segnata da un carattere speciale (EOF), ma 
si rileva in base all'esito delle operazioni di lettura.

\begin{lstlisting}[style=cstyle]
// legge n sequenze di byte di lunghezza Dato
fread(indirizzo, sizeof(Dato), n, fp); 

// scrive n sequenze di byte di lunghezza Dato
fwrite(indirizzo, sizeof(Dato), n, fp);
\end{lstlisting}

la funzione \texttt{fwrite} restituisce il numero di elementi scritti con successo (che possono essere meno di n),
prelevando i byte a partire dall'indirizzo specificato.
Analogamente, \texttt{fread} restituisce il numero di elementi letti con successo, che possono
essere meno di n se si raggiunge la fine del file. Controllare il valore restituito è 
il solo modo per sapere cosa è stato letto, in particolare per capire se il file è terminato.



\subsection{Accesso Diretto: fseek e ftell}

\begin{lstlisting}[style=cstyle]
fseek(fp, 0L, SEEK_END);   // va alla fine del file
pos = ftell(fp);           // posizione corrente
rewind(fp);                // equivalente a fseek(fp,0,SEEK_SET)
\end{lstlisting}

\textbf{Nota:} La funzione \texttt{fseek} permette di spostare il puntatore di lettura/scrittura 
in una posizione arbitraria del file. La funzione fornisce tre punti di partenza per lo spostamento:
\begin{itemize}
    \item \texttt{SEEK\_SET}: dall'inizio del file.
    \item \texttt{SEEK\_CUR}: dalla posizione corrente.
    \item \texttt{SEEK\_END}: dalla fine del file.
\end{itemize} 

\section{Argomenti da Riga di Comando}

È possibile passare informazioni a un programma direttamente da terminale:

\begin{lstlisting}[style=cstyle]
./programma arg1 arg2 arg3
\end{lstlisting}

La funzione \texttt{main} può ricevere due parametri:

\begin{lstlisting}[style=cstyle]
int main(int argc, char *argv[]) {
    ...
}
\end{lstlisting}

\begin{itemize}
    \item \texttt{argc} = numero totale di argomenti (il primo è il nome del programma);
    \item \texttt{argv} = array di stringhe con gli argomenti.
\end{itemize}

\textbf{Esempio}

\begin{lstlisting}[style=console]
    ./programma prova.txt
\end{lstlisting}

\begin{lstlisting}[style=cstyle]
int main(int argc, char *argv[]) {
    if (argc != 2) {
        printf("Uso: %s nomefile\n", argv[0]);
        return 1;
    }

    FILE *fp = fopen(argv[1], "r");
    if (fp == NULL) {
        printf("Errore apertura file!\n");
        return 1;
    }

    /* elaborazione */
    fclose(fp);
    return 0;
}
\end{lstlisting}

\textbf{Conversione di argomenti numerici}

\begin{lstlisting}[style=cstyle]
int n = atoi(argv[1]);
float x = atof(argv[2]);
char c = argv[3][0];
\end{lstlisting}

\section{Variabili Locali, Statiche e Globali}

Il linguaggio C distingue diverse categorie in base allo \emph{tempo di vita} e allo \emph{scope}.

\subsection{Variabili Locali Automatiche}

Vivono solo all'interno della funzione in cui sono dichiarate.
Sono allocate sullo stack.

\begin{lstlisting}[style=cstyle]
void f() {
    int x = 0; // nuova x a ogni chiamata
}
\end{lstlisting}

\subsection{Variabili Locali Statiche}

Mantengono il proprio valore anche dopo la fine della funzione.

\begin{lstlisting}[style=cstyle]
void f() {
    static int x = 0;
    x++;
    printf("%d\n", x);
}
\end{lstlisting}

Ogni chiamata stampa un valore crescente.

\subsection{Variabili Globali}

Sono dichiarate fuori da qualsiasi funzione.
Sono visibili da tutto il file (scope di file).

\begin{lstlisting}[style=cstyle]
int contatore = 0; // globale

void f() { contatore++; }
\end{lstlisting}

\paragraph{Attenzione}  
L’uso indiscriminato di variabili globali rende difficile il debug e la manutenibilità del codice.

\subsection{Blocchi di Istruzioni e Regole dello Scope}

Un \textbf{blocco} è un insieme di istruzioni racchiuse tra \texttt{\{} e \texttt{\}}.

Le variabili dichiarate in un blocco esistono solo all’interno di quel blocco.

\begin{lstlisting}[style=cstyle]
int x = 1;
if (x > 0) {
    int x = 5;   // nasconde la x esterna
    printf("%d\n", x); // stampa 5
}
printf("%d\n", x);     // stampa 1
\end{lstlisting}

Una dichiarazione interna con lo stesso nome di una variabile esterna \textbf{nasconde} temporaneamente quella esterna.

\section{Esercizi}

% Facili
\begin{Exercise}[label={ex:lettura_scrittura}]
    Si scriva un programma che, aperto un file di testo di nome "output.txt", in 
    modalità scrittura, provveda a leggere da input delle parole separate da spazi
    (stringhe di al più 63 caratteri) e le scriva nel file di testo.
    La lettura termina quando l'utente inserisce la parola "fine".
    Successivamente, il programma riapre il file in modalità lettura e stampa a video
    tutte le parole lette dal file.
\end{Exercise}

\begin{Answer}[ref={ex:lettura_scrittura}]
  	\lstinputlisting[style=cstylesmall]{src/lettura_scrittura.c}
\end{Answer}
\begin{Exercise}[label={ex:massimo}]
    Si scriva un programma che apre un file di testo il cui
    nome è input1.txt contenente una sequenza di lunghezza indefinita (0 o più)
    di numeri interi. Il programma individua nel file il valore massimo e lo scrive in
    un altro file output1.txt insieme al numero di volte che vi compare.
\end{Exercise}

\begin{Answer}[ref={ex:massimo}]
  	\lstinputlisting[style=cstylesmall]{src/massimo.c}
\end{Answer}
\begin{Exercise}[label={ex:matrice}]
    Una matrice sparsa è una matrice in cui la maggior parte degli elementi è costituita da zeri. 
    Spesso una matrice sparsa viene rappresentata in formato COO (Coordinate List), 
    ovvero come una sequenza di triple (i, j, valore) che indicano che, nella posizione (i, j), 
    è presente un valore diverso da zero (valore $\neq$ 0). Si consideri un file binario denominato 
    "input6.bin", che contiene due interi rappresentanti il numero di righe e di colonne della matrice, 
    seguiti dalla matrice stessa rappresentata in modalità ordinaria (cioè come sequenza di r * c elementi,
    dove r è il numero di righe e c il numero di colonne), i cui elementi sono numeri
    in virgola mobile a doppia precisione.
    Si scriva un programma che legga tale file e salvi in un secondo file di testo denominato "coo6.txt" 
    la rappresentazione COO della matrice.
\end{Exercise}

\begin{Answer}[ref={ex:matrice}]
  	\lstinputlisting[style=cstylesmall]{src/matrice.c}
\end{Answer}


\begin{Exercise}[label={ex:studenti}]
    Dato un file di testo "dati.txt" contenente i dati relativi agli studenti
    immatricolati al primo anno della Facoltà di Ingegneria. In particolare, le informazioni
    sono memorizzate nel file con uno studente per riga nel seguente formato:
    \begin{itemize}
        \item matricola (intero);
        \item corso di laurea (intero);
    \end{itemize}
    Sia dato un secondo file di testo, "indirizzi.txt" che contiene l'indirizzo di ogni studente.
    In particolare, in ogni linea, separati da uno spazio: 
    \begin{itemize}
        \item matricola (intero);   
        \item nome (stringa di al più 20 caratteri);
        \item cognome (stringa di al più 30 caratteri);
        \item via (stringa di al più 30 caratteri);
        \item città (stringa di al più 20 caratteri);
        \item cap (intero).
    \end{itemize}
    Si scriva un programma che:
    \begin{enumerate}
        \item a partire dai file "dati.txt" e "indirizzi.txt" costruisca una tabella
        contenente per ogni studente: matricola, nome, cognome, via, città, cap e corso di laurea.
        \item a partire dalla tabella e dato input un intero C che rappresenta un corso di laurea,
        stampi la percentuale di studenti iscritti a quel corso.
        \item scriva su un terzo file di testo "milano.txt", nome, cognome e numero di matricola
        di tutti gli studenti che abitano a Milano.
    \end{enumerate}
\end{Exercise}

\begin{Answer}[ref={ex:studenti}]
  	\lstinputlisting[style=cstylesmall]{src/studenti.c}
\end{Answer}
\begin{Exercise}[label={ex:telefoni}]
    Una società di telefonia cellulare gestisce un programma di premiazione per "utenti fedeli". In particolare, per
    ogni cliente viene salvato su un file binario "punti.dat" il nome del cliente (al massimo 31 caratteri) e un
    numero intero che rappresenta i punti accumulati. Tali informazioni sono organizzate come una struttura \texttt{User}, opportunamente definita.

    Si scriva una funzione \texttt{readPoints()} che ricevuto in ingresso il nome di un file \texttt{usersFile}, un array \texttt{results} di
    strutture \texttt{User}, la dimensione massima dell’array \texttt{maxDim}, e un limite
    inferiore di punti \texttt{minPoints}, copi nell’array \texttt{results} i dati dei clienti che hanno almeno i punti specificati da \texttt{minPoints}.

    La funzione deve restituire come risultato il numero di utenti con almeno
    \texttt{minPoints}. Si noti che tale risultato rappresenta anche la dimensione
    logica dell’array results. Qualora il file non sia accessibile, la funzione deve restituire il valore -1.

    Si scriva poi un programma \texttt{main()} che chieda all’utente il numero
    di clienti salvati sul file (tale numero sarà noto solo a tempo di
    esecuzione), e allochi dinamicamente un vettore V di \texttt{User}
    sufficientemente grande per poter contenere, nel caso peggiore, i
    dati di tutti gli utenti salvati in \texttt{usersFile}. Il programma dovrà poi
    chiedere all’utente il minimo punteggio e, utilizzando la funzione
    \texttt{readPoints()}, leggere da file e memorizzare in V i dati degli
    utenti che hanno almeno il punteggio minimo specificato. Il
    programma infine deve stampare a video il nome ed il punteggio
    degli utenti contenuti in V se e solo se il nome comincia per “Me”.

\end{Exercise}

\begin{Answer}[ref={ex:telefoni}]
  	\lstinputlisting[style=cstylesmall]{src/telefoni.c}
\end{Answer}
\begin{Exercise}[label={ex:cifrario}]
    Si scriva un programma che decifra o cifra un file di testo applicando uno scorrimento ai caratteri alfabetici.
    Il programma deve accettare i parametri \emph{in qualsiasi ordine} tramite dei flag da linea di comando:
    \begin{itemize}
        \item \texttt{-i <inputfile>} : specifica il file di input
        \item \texttt{-o <outputfile>} : specifica il file di output
        \item \texttt{-k <shift>} : specifica la chiave
        \item \texttt{-d} : decifra il file (se non presente, cifra)
    \end{itemize}

    \textbf{Esempio:}
    \begin{lstlisting}[style=console]
        # Cifra il file msg.txt con chiave 3 e salva il risultato in cif.txt
        ./programma -i msg.txt -o cif.txt -k 3
        
        # Decifra il file cif.txt con chiave 5 e salva il risultato in msg.txt
        ./programma -k 5 -d -o msg.txt -i cif.txt
    \end{lstlisting}
\end{Exercise}

\begin{Answer}[ref={ex:cifrario}]
  	\lstinputlisting[style=cstylesmall]{src/cifrario.c}
\end{Answer}
\begin{Exercise}[label={ex:voti}]
    Si scriva un programma che gestisca un database di voti universitari salvati su un file "registro.txt".
    Il programma deve gestire due modilità operative distinte:
    \begin{itemize}
        \item Modalità aggiunta (-add): aggiunge un nuovo studente al registro. Richiede tre parametri: nome, cognome e voto. 
        \item Modalità ricerca (-find): cerca e stampa a video tutti gli studenti che hanno un voto superiore o uguale ad una soglia specificata.
    \end{itemize}

    Si scriva una funzione \texttt{aggiungi\_studente()} che aggiunge un nuovo studente al file "registro.txt" 
    e una funzione \texttt{cerca\_studenti()} che apre e cerca all'interno di "registro.txt" 
    tutti gli studenti con voto superiore o uguale alla soglia specificata. Si gestiscano le due modalità come
    argomenti da riga di comando e si aggiunga la possibilità di utilizzare un file di registro personalizzato usando come flag "-f".

    \textbf{Esempio:}
    \begin{lstlisting}[style=console]
        # Usa il file di default, aggiunge Mario Rossi con 28
        ./programma -add Mario Rossi 28

        # Usa un file diverso, cerca voti >= 25
        ./programma -f esame_c.txt -find 25
    \end{lstlisting}
    
\end{Exercise}

\begin{Answer}[ref={ex:voti}]
  	\lstinputlisting[style=cstylesmall]{src/voti.c}
\end{Answer}
\begin{Exercise}[label={ex:archivio}]
    Un file di testo "archivio.txt" contiene i dati (primo autore, titolo, numero di copie possedute, 
    numero di copie in prestito) relativi ai differenti volumi conservati presso una biblioteca.
    Più precisamente, ogni riga del file contiene nell'ordine, separati da uno spazio bianco:
    \begin{itemize}
        \item autore (al più 23 caratteri senza spazi)
        \item titolo (al più 55 caratteri senza spazi)
        \item numero di copie possedute (intero)
        \item numero di copie in prestito (intero)
    \end{itemize}

    Si realizzi un programma C che:
    \begin{enumerate}
        \item Legga il contenuto di "archivio.txt" e costruisca in memoria un vettore V di strutture
        corrispondenti (si supponga che il file "archivio.txt" non possa contenere più di 30 righe). 
        Si stampi a video il contenuto del vettore.
        \item A partire da V, costruisca una lista L di interi contenente per ciascun volume il 
        numero di copie disponibili nella biblioteca, ovvero la differenza fra il numero di copie 
        possedute e il numero di copie in prestito. Si stampi a video il contenuto della lista L.
        \item Utilizzando L per ottenere la somma delle copie disponibili e V per la somma delle
        copie possedute, calcoli il rapporto fra volumi disponibili e volumi posseduti.
        \item  Utilizzando la lista di interi L, stampi il numero di riga di "archivio.txt" 
        relativo al volume con più copie disponibili. In caso di più volumi con pari numero di copie 
        disponibili, qualunque riga relativa a questi ultimi è considerata una risposta corretta.
    \end{enumerate}
\end{Exercise}

\begin{Answer}[ref={ex:archivio}]
  	\lstinputlisting[style=cstylesmall]{src/archivio.c}
\end{Answer}

% Difficili


\clearpage
\section*{Soluzioni}
\shipoutAnswer


\end{document}